\begin{sitedescription}{EP}

\begin{center}
\includegraphics[height=3cm]{Participants/Logos/EP.png}
\end{center}

\'Ecole Polytechnique (l'X), is the leading member of Grandes \'Ecoles in France
for science and technology according to all French rankings last year, ranked
2nd best small university in the world by Times Higher Education in 2018, and
16th best European university by QS World University Rankings in 2018. With
fewer than 3000 Bachelor, Masters, and PhD students enrolled at any time, the
institution has produced 4 Fields Medalists, 4 Nobel prize winners, and 3 French
presidents, ranking it 6th worldwide in terms of number of Nobel prize
recipients. L'X is composed of 22 laboratories supporting research in physics,
engineering, mathematics, biology, and chemistry (among others) and is connected
to several French national research institutions such as CNRS, CEA, and INRIA.
L'X prides itself on promoting multidisciplinary research through collaboration
between its many labs and with external, national and international partners and
a strong connection with industry and start-up.

The Mathematics department at Ecole polytechnique has started a reform of their
various teaching offer based on Jupyter. Many courses from the
Bachelor program, Engineering school, and the Master program have
already begun relying on a strong use of Jupyter notebooks /
JupyterHub\footnote{\url{http://www.cmap.polytechnique.fr/~massot/Personal_web_page_of_Marc_Massot/MAP551}} and this will continue with a strong support of the Dean of
undergraduate studies and of graduate studies. In addition, several software and
research engineers in applied mathematics have been recruited and
participate, together with a system engineer, to this effort in order to help in terms of
building an infrastructure dedicated to Jupyter.
This infrastructure together with support from the head of the Ecole polytechnique,
fosters the dissemination into various other
departments (Physics, Mechanical Engineering, Biology...), where some
Jupyter-based courses are already starting.
A community is emerging, notably through the strong engagement of
students in the computer science club of \'Ecole polytechnique (Binet R\'eseau).

\subsubsection*{Curriculum vitae}
% Curriculum of the personnel at this institution. This includes
% to-be-hired people for which there is a tentative candidate.

\begin{participant}[type=leadPI,PM=8,gender=male]{Loic Gouarin}

  Loic Gouarin is Research Engineer in scientific computing at CMAP (Centre de
  Mathématiques Appliquées) which is a Joint Research Unit
  (UMR7641) of the CNRS and \'Ecole polytechnique. He works on several
  scientific computing open-source projects in different fields such as
  Lattice-Boltzmann methods, Stokes solvers for fluid particles interaction,
  adaptive mesh refinement, ...

  He is also director of the ``GdR Calcul'' where his role is to animate the
  scientific and high performance computing community in France, in particular
  by organising conferences, meetings, and seminars. In this context, he
  organises himself 3 to 4 training and development workshops per year, and
  promotes the use of Python and c++ for teaching and research in France.
  
  For several years, he has been very involved in promoting the Jupyter project
  and its use for teaching and research in the French community. He is one of
  the core developers of xeus-cling and he is working on the possibility of
  easily deploying a JupyterHub or BinderHub on academic clouds. He also
  believes that reproducible research is an essential part of promoting new
  computation codes, new numerical methods, ... introduced in related
  publications and therefore he uses Jupyter as a first approach to achieve this.
  
  \end{participant}
  
  %%% Local Variables:
  %%% mode: latex
  %%% TeX-master: "../proposal"
  %%% End:
  
\begin{participant}[type=R,PM=4,gender=male]{David Delavennat}

    David Delavennat is Research Engineer in scientific infrastructures at CMLS
    (Centre de Mathématiques Laurent Schwartz) which is a Joint Research Unit
    (UMR7640) of the CNRS and \'Ecole polytechnique. He
    works on several french scientific infrastructures projects using
    microservices, virtualization, container and cloud technologies.

    He is coordinator of ARGOS, the Paris-South business network targeted to the
    academic System and Network Administrators. He animates this community,
    organising conferences and meetings He represents ARGOS at RESINFO, the
    french business network of regional SNA business networks. In this context,
    he participate each year to the organization of National Teaching Actions.

    For several years, he has been very involved in promoting the DevOps
    paradigm and Continuous Integration for deploying production infrastructures
    on academic Openstack and Kubernetes environment for the French Mathematical
    community. He is working on simplifying the deployment of a JupyterHub or
    BinderHub onto academic PaaS and Iaas. He also believes that reproducible
    infrastructures is an essential part of reproducible science.

\end{participant}
\input{CVs/Marc.Massot.tex}
\begin{participant}[type=R,PM=2,gender=male]{Laurent Series}

    Laurent Series is Research Engineer in scientific computing at CMAP (Centre 
    de Mathématiques Appliquées) which is a Joint Research Unit
    (UMR7641) of the CNRS and \'Ecole polytechnique. He works on several
    scientific computing open-source projects in different fields like
    finite element method for solid mechanics, multiresolution for 
    reaction-diffusion equations modeling multi-scale reaction wave ...
  
    He was technical responsible of Computing Center (M\'esocentre) of Ecole
    CentraleSup\'elec from 2009 to 2018. In this context, he organises and 
    participates in the user support team (assistance in porting codes, 
    parallelization, vectorisation, optimisation). Since his arrival at \'Ecole 
    Polytechnique, he is involved in the process of creation of a new Computing 
    Center.
    
    For several years, he has been very involved in promoting the Jupyter project
    by writing Jupyter notebooks for course (such as "Dynamical systems for the 
    modeling and simulation of multi-scale reacting media" MAP551) and by promoting, 
    among researchers, its use to share their work. He is also involved with the 
    direction of \'Ecole polytechnique for the project of building an infrastructure 
    for Jupyter at the level of the school, for its use in both research and teaching.
  
\end{participant}

% For other to-be-hired person, please include here something like:
\begin{participant}[type=R,PM=18,salary=5500]{DevOps Engineer}
  We intend to hire a DevOps engineer for 18 months during the project to be supervised at \'Ecole polytechnique to work on developments defined in Tasks \taskref{core}{jh-bh-conv} and \taskref{eosc}{jh-bh-deployment}.
\end{participant}

\begin{participant}[type=R,PM=18,salary=5500]{Software Engineer}
  We intend to hire a DevOps engineer for 18 months during the project to be supervised at \'Ecole polytechnique to work on developments defined in Task \taskref{ecosystem}{teaching-tools}.
\end{participant}

\subsubsection*{Publications, products, achievements}

\begin{compactenum}
\item D. Delavennat, L. Gouarin, G. Philippon, \emph{Deploying JupyterHub with Kubernetes on OpenStack} \newline
\url{https://blog.jupyter.org/how-to-deploy-jupyterhub-with-kubernetes-on-openstack-f8f6120d4b1}
\item JupyterDay at \'Ecole polytechnique in 2018 \newline
\url{http://www.cmap.polytechnique.fr/~massot/Personal_web_page_of_Marc_Massot/JupyterX.html}
\item L. Gouarin, \emph{C++ course using xeus-cling} \newline
\url{https://github.com/gouarin/cours_cpp_moderne}
\item M.Massot, L. Series, \emph{Systèmes dynamiques pour la modélisation et la simulation des "milieux réactifs" multi-échelles} \newline
\url{http://www.cmap.polytechnique.fr/~massot/Personal_web_page_of_Marc_Massot/MAP551}
\end{compactenum}

\subsubsection*{Relevant projects or activities}

\begin{compactenum}
\item Through a previous position at Paris Sud, Loïc Gouarin is a
  participant of OpenDreamKit (GA No. 676541) Open Digital Research
  Environment Toolkit for the Advancement of Mathematics.
\end{compactenum}

\subsubsection*{Significant infrastructure}

\site{EP} has strongly contributed to the deployment of JupyterHub on
\site{UPSUD}'s local OpenStack based cloud infrastructure
\software{Cloud@VD}. \site{EP} invested last year in the purchase of 100 cores for this infrastructure in order to study how to deploy JupyterHub and BinderHub and start to offer to researchers and students this kind of service.
In 2019, \site{EP} will have its own OpenShift cloud infrastructure with 200-300 cores dedicated to teaching with Jupyter.

\end{sitedescription}
%%% Local Variables:
%%% mode: latex
%%% TeX-master: "../proposal"
%%% End:
